\section{Calculadora do Cidadão}
\label{cap-solucoes-calculadora}

A Calculadora do Cidadão é uma ferramenta disponibilizada pelo Banco Central do Brasil para simular situações do cotidiano financeiro. Na modalidade de financiamento com prestações fixas, a ferramenta permite relacionar o valor financiado, a taxa mensal de juros, o número de meses e o valor da prestação. O usuário informa os dados conhecidos para calcular a variável de interesse \cite{bcbCalculadora}.


Para este projeto, a ferramenta serve como referência para a apresentação das informações de uma simulação de crédito. A identificação da taxa mensal, do prazo e do valor financiado ajuda o usuário a reconhecer quais dados foram utilizados para obter o resultado.

\section{Comparação com a proposta do simulador bancário}
\label{cap-solucoes-comparacao}

A Calculadora do Cidadão oferece uma referência para estruturar a interação entre os dados informados e o resultado de um cálculo financeiro \cite{bcbCalculadora}. O simulador proposto neste trabalho situa esse tipo de cálculo em um ambiente didático que também representa contas, saldos e movimentações fictícias.


A aplicação prevê dois tipos de atividade. Na primeira, o estudante realiza pagamentos e transferências e verifica seus efeitos sobre o saldo e o extrato. Na segunda, altera parâmetros de uma simulação de crédito e compara os resultados obtidos. Essas atividades permitem trabalhar tanto a interpretação de movimentações quanto a análise de taxa, prazo e custo.